\documentclass[11pt]{article}

\usepackage[margin=1in]{geometry}

\usepackage{amsmath,amssymb,amsthm}
\usepackage{booktabs}
\usepackage{array}
\usepackage[table]{xcolor}
\usepackage{siunitx}
\usepackage{enumitem}
\usepackage[expansion=false,protrusion=true]{microtype}
\DisableLigatures[f]{encoding = *}
\usepackage[hidelinks]{hyperref}
\usepackage{titlesec}
\usepackage{caption}

% --- F1R3FLY brand accent ---
\definecolor{f1r3red}{HTML}{D7263D}
\definecolor{f1r3ink}{HTML}{1A1A1A}
\definecolor{f1r3ash}{HTML}{6B6B6B}
\definecolor{coolrow}{HTML}{EAF5EC}
\definecolor{hotrow}{HTML}{FBE9EC}

\titleformat{\section}{\color{f1r3red}\large\bfseries}{\thesection}{0.6em}{}
\titleformat{\subsection}{\color{f1r3ink}\normalsize\bfseries}{\thesubsection}{0.6em}{}

\hypersetup{colorlinks=true,linkcolor=f1r3red,citecolor=f1r3red,urlcolor=f1r3red}

\newcommand{\shards}{S}
\newcommand{\agentcost}{a}
\sisetup{group-separator={,},group-minimum-digits=4}

\newtheorem{remark}{Remark}

\title{\color{f1r3red}\bfseries Managing the Shard Tree\\[2pt]
\large\color{f1r3ink}\mdseries Agents, DevOps, and the automation crossover ---
a management-economics companion to occupancy-driven sharding in F1R3FLY}
\author{\normalsize F1R3FLY Industries --- working note}
\date{\normalsize \today}

\begin{document}
\maketitle

\noindent\textit{A short note, and a companion to ``Rent and Shard Splitting.'' That note priced
\emph{running} a shard tree; this one prices \emph{operating} it---who provisions the nodes, watches
the tables fill, fires the split, migrates the cold tenants, and answers the pager at 3 a.m. The
arithmetic in \S\ref{sec:cross} is an engineering estimate built on public 2026 figures (Anthropic
API list prices, OCI GPU list prices, and US SRE/DevOps compensation surveys); every assumption is a
lever in the companion workbook so the numbers can be re-derived or re-priced.}

\section{The problem in one paragraph}

A single F1R3FLY shard is ten nodes costing $C\approx\$\num{4160}$/month
(``Rent and Shard Splitting,'' \S5). But a shard is never alone for long: as occupancy ratchets
toward the split threshold $\alpha^\star$, the tree fans out, and the sector roadmap in the
sharding-economics workbook carries the AI engine alone from one shard to $\sim\!\num{334}$ by
year five, with the cross-shard \emph{connector} layer reaching $\sim\!\num{12000}$ shards at
maturity. Each of those shards must be stood up, monitored, split, and---per the rent
lifecycle---have its tenants demoted and evicted on schedule. That operational work can be done by
a human DevOps team, by an autonomous agent, or by a mix. The three differ not in \emph{what} they
do but in \emph{how their cost scales}: human cost is a step function with a high fixed floor and a
steep per-shard slope; agent cost is nearly flat, a fixed platform charge plus pennies of inference
per shard. Where those curves cross is the whole question, and it lands at a surprisingly small
tree.

\section{Three regimes}
\label{sec:regimes}

We cost three ways of operating a tree of $\shards$ homogeneous shards.

\begin{description}[leftmargin=0pt,labelsep=0.6em]
  \item[\textnormal{\textbf{(1) Full automation.}}] An agent provisions hardware, stands up nodes,
  and manages the whole tree---health, splits, rent decisions, tenant migration. No in-house
  operators; physical failures are absorbed by the cloud provider's own hardware replacement, and a
  thin premium-support retainer covers SLA disputes. This is the cloud-native limit. It is a
  \emph{price floor}, not a deployable plan, because an agent cannot currently re-seat a failed DIMM
  or litigate a credit with a carrier; we cost it anyway, because it is the asymptote the other two
  are measured against.
  \item[\textnormal{\textbf{(2) Agent in the mix (hybrid).}}] The agent does the routine---all
  monitoring, provisioning, split and migration decisions---and escalates the exceptions to a thin
  human team sized by the \emph{exception rate}, not by the shard count. This is the realistic
  deployment.
  \item[\textnormal{\textbf{(3) DevOps team.}}] Humans provision and manage the tree with
  conventional infrastructure-as-code. The agent platform is not built; cost is people.
\end{description}

\section{The cost models}
\label{sec:models}

Let $W$ be the fully-loaded monthly cost of one SRE, $\sigma$ a span of control (shards one engineer
can carry), $\varphi$ a coverage floor (minimum head\-count for a 24/7 rotation), $P$ the fixed cost
of \emph{having} an ops agent (the engineers who build and maintain the harness, evals, and
guardrails, plus telemetry), $\rho$ a premium-support retainer, and $\agentcost$ the agent's
inference cost to manage one shard for one month. Then

\begin{align}
\text{DevOps:}     \quad M_d(\shards) &= W\cdot\max\!\bigl(\varphi_d,\ \lceil \shards/\sigma_d\rceil\bigr), \label{eq:devops}\\
\text{Full-auto:}  \quad M_f(\shards) &= P + \rho + \agentcost\,\shards, \label{eq:full}\\
\text{Hybrid:}     \quad M_h(\shards) &= P + \agentcost\,\shards + W\cdot\max\!\bigl(\varphi_h,\ \lceil \shards/\sigma_h\rceil\bigr). \label{eq:hybrid}
\end{align}

The structure is the entire story. Equation~\eqref{eq:devops} is flat at $W\varphi_d$ until the tree
outgrows the on-call floor, then climbs one engineer per $\sigma_d$ shards---marginal cost
$W/\sigma_d$ per shard. Equation~\eqref{eq:full} is a fixed intercept plus a gentle slope
$\agentcost$. Equation~\eqref{eq:hybrid} pays \emph{both} a platform intercept and a human term, but
the human term uses $\sigma_h\gg\sigma_d$ because the agent has absorbed everything routine: one
operator now backstops hundreds of shards instead of dozens.

\subsection{What the agent's \texorpdfstring{$\agentcost$}{a} actually is}
The per-shard inference cost is a transparent token sum over the agent's activities $k$ (routine
monitoring, hard decisions, provisioning), each with $n_k$ calls per shard-month, input/output token
counts $\tau^{\text{in}}_k,\tau^{\text{out}}_k$, a cached-input fraction $\beta_k$, and model prices
$p^{\text{in}}_k,p^{\text{out}}_k$ per million tokens with cache multiplier $\gamma$ (prompt caching
bills cached input at \SI{10}{\percent}):
\[
\agentcost \;=\; \sum_k n_k\,\frac{\bigl[\beta_k\gamma + (1-\beta_k)\bigr]\,\tau^{\text{in}}_k\,p^{\text{in}}_k
\;+\; \tau^{\text{out}}_k\,p^{\text{out}}_k}{10^6}.
\]
With routine monitoring on Haiku 4.5, hard decisions (split / migrate / triage) on Opus 4.8, and
provisioning on Sonnet 4.6 at June 2026 list prices~\cite{anthropicprice2026}, the workbook returns

\begin{table}[h]
\centering
\caption*{\textbf{Table 1.} Agent inference cost to manage one shard for one month.}
\begin{tabular}{@{}lrrr@{}}
\toprule
Activity & Calls/shard-mo & Model & Cost (\$/shard-mo) \\
\midrule
Routine monitoring   & \num{720} & Haiku 4.5  & \$2.61 \\
Hard decisions       & \num{6}   & Opus 4.8   & \$1.40 \\
Provisioning / stand-up & \num{0.5} & Sonnet 4.6 & \$0.17 \\
\midrule
\textbf{Total $\agentcost$} & & & \textbf{\$4.17} \\
\bottomrule
\end{tabular}
\end{table}

\noindent Four dollars and change. Against an infrastructure bill of $C\approx\$\num{4160}$ per
shard, the agent's management overhead is \emph{one tenth of one percent}. The same shard managed at
the margin by a human costs $W/\sigma_d\approx\$\num{667}$/month---about \SI{16}{\percent} of infra,
and roughly \num{160} times the agent. That ratio, not any subtlety of the curves, is why the
crossover exists and why it is early.

\section{The crossover}
\label{sec:cross}

\paragraph{Parameters.} Loaded SRE $W=\$\num{20000}$/month (a \$\num{150000} total-comp engineer at a
$1.6\times$ loading factor~\cite{builtin2026,kore2026}); DevOps span $\sigma_d=30$ shards/SRE with a
lean 24/7 floor $\varphi_d=2$; hybrid span $\sigma_h=300$ with exception floor $\varphi_h=1$; agent
platform $P=\$\num{28000}$/month ($1.25$ platform FTEs plus \$\num{3000} tooling); retainer
$\rho=\$\num{2000}$. All are levers; \S\ref{sec:sens} sweeps the ones that move the answer.

\begin{table}[h]
\centering
\caption*{\textbf{Table 2.} Monthly cost to manage the whole tree. ``Feasible best'' chooses between
the two deployable regimes (DevOps, Hybrid); Full-auto is shown as the floor.}
\setlength{\tabcolsep}{6pt}
\renewcommand{\arraystretch}{1.12}
\begin{tabular}{@{}rrrrl@{}}
\toprule
Shards $\shards$ & DevOps & Full-auto (floor) & Hybrid & Feasible best \\
\midrule
\num{1}     & \$\num{40000}      & \$\num{30004}  & \$\num{48004}  & DevOps \\
\num{25}    & \$\num{40000}      & \$\num{30104}  & \$\num{48104}  & DevOps \\
\num{50}    & \$\num{40000}      & \$\num{30208}  & \$\num{48208}  & DevOps \\
\num{60}    & \$\num{40000}      & \$\num{30250}  & \$\num{48250}  & DevOps \\
\rowcolor{coolrow}\num{90} & \$\num{60000} & \$\num{30375} & \$\num{48375} & Hybrid \\
\num{100}   & \$\num{80000}      & \$\num{30417}  & \$\num{48417}  & Hybrid \\
\num{250}   & \$\num{180000}     & \$\num{31042}  & \$\num{49042}  & Hybrid \\
\num{334}   & \$\num{240000}     & \$\num{31392}  & \$\num{69392}  & Hybrid \\
\num{1000}  & \$\num{680000}     & \$\num{34166}  & \$\num{112166} & Hybrid \\
\num{5000}  & \$\num{3340000}    & \$\num{50832}  & \$\num{388832} & Hybrid \\
\rowcolor{hotrow}\num{12074} & \$\num{8060000} & \$\num{80305} & \$\num{898305} & Hybrid \\
\bottomrule
\end{tabular}
\end{table}

\paragraph{Where it crosses.} Below $\sim\!60$ shards a lean human team is the cheapest
\emph{deployable} option: the tree is small enough that two on-call engineers cover it, and that
\$\num{40000} undercuts the \$\num{48000} floor of standing up \emph{and} owning an agent platform.
The flip comes the moment DevOps must hire its third engineer---at $\shards>2\sigma_d=60$---after
which human cost climbs a step every thirty shards while hybrid barely moves. Dropping the ceilings
and floors, the crossover has a clean closed form:
\[
\shards^\star \;\approx\; \frac{P+W}{\,W/\sigma_d - \agentcost\,}
\;=\; \frac{\num{28000}+\num{20000}}{\,\num{20000}/30 - 4.17\,}
\;\approx\; \num{72}\ \text{shards}.
\]
The denominator is the per-shard \emph{advantage} of an agent over a human ($\$\num{667}-\$4$); the
numerator is the fixed price of admission to automation. Both sides are small numbers, so the
quotient is small: F1R3FLY crosses into agent-favourable territory while its tree is still in the
low dozens---well inside the year-four AI ramp, and an order of magnitude before the connector layer
exists at all.

\paragraph{The two asymptotes.} At the high end the regimes separate brutally. At the connector-scale
$\num{12074}$-shard tree, a pure DevOps organisation needs \num{403} engineers and
\$\num{8.06}M/month; the hybrid model needs \num{41} (one platform team plus exception coverage) and
\$\num{0.90}M---\SI{11}{\percent} of the DevOps bill; full automation, were it feasible alone, would
cost \$\num{80000}, or \SI{1}{\percent}. As a fraction of the \$\num{50}M infrastructure spend at
that scale, human-bound management is a persistent \SI{16}{\percent} tax, while agent-led management
falls below \SI{2}{\percent} and keeps falling. Management is a major line item exactly when it is
human and the tree is large; it is a rounding error exactly when it is automated.

\begin{remark}[The agent is already in the design]
The hard decisions priced on Opus above are not generic ``ops.'' They are the very events the rent
companion specifies: \emph{when to split} (the $\alpha^\star$ test), \emph{which tenants to demote}
and \emph{how far} (the bid-rent sort), and \emph{when to evict} (the depletion cascade). These are
local, state-dependent, content-addressed decisions over a tuplespace---precisely the shape a
reactive agent handles well, and precisely the shape that scales badly across a human on-call
rotation. The split threshold and the management model are the same artefact seen twice: occupancy
fires the split; the agent executes it; rent prices the migration. Automating the tree is not a new
system bolted on, but the reducer's own lifecycle given an operator.
\end{remark}

\section{Why hybrid, and what stays human}
\label{sec:hybrid}

Full automation wins on dollars at every scale (Table 2, floor column) for the same reason it is not
deployable: it books no humans. The residual that keeps humans in the loop is small but irreducible,
and it is physical or adversarial rather than computational:
\begin{itemize}[leftmargin=1.4em]
  \item \textbf{Hardware the provider won't touch.} On managed cloud the carrier swaps failed
  disks and hosts; on owned or co-located metal someone must. This is the difference between
  $\rho$ (a support line item) and an on-site hands contract.
  \item \textbf{SLA and contract disputes.} Negotiating a credit for a region outage, or a
  commitment discount, is a human-to-human commercial act. The agent can assemble the evidence; a
  person signs.
  \item \textbf{Novel incidents and change authority.} The first occurrence of a failure mode, and
  any irreversible action (key rotation, consensus reconfiguration), wants a human in the approval
  path until the agent's eval coverage earns autonomy.
\end{itemize}
Hybrid prices exactly this: it pays the agent's flat $\agentcost\shards$ for everything routine and
keeps a team sized by how often the above fires---captured here as the tenfold span
$\sigma_h=300$. Raise the exception rate (owned hardware, regulated change control, immature evals)
and $\sigma_h$ falls toward $\sigma_d$, pulling hybrid back toward DevOps; lower it (fully managed
cloud, mature autonomy) and $\sigma_h$ rises, pulling hybrid toward the full-auto floor. The single
parameter $\sigma_h$ is the dial between ``agent with human backup'' and ``humans with agent help.''

\section{Sensitivity: the levers that move the answer}
\label{sec:sens}

\paragraph{The DevOps floor.} The crossover sits at $\sim\!60$ shards \emph{because} we modelled a
lean two-person rotation. A team holding a strict 24/7 SLA at $\varphi_d=4$ pays \$\num{80000} from
the first shard, and hybrid wins essentially everywhere; conversely a single-engineer hobby
deployment ($\varphi_d=1$) pushes the crossover out past \num{100} shards. The floor is the most
consequential single assumption, and it is a policy choice (what coverage you promise), not a
technical one.

\paragraph{The agent platform $P$.} Halving the platform cost (a shared agent across many trees,
amortised tooling) moves $\shards^\star$ down proportionally and makes hybrid competitive against
even a one-person team. Doubling it (bespoke harness, heavy eval burden) roughly doubles the
crossover. $P$ is a one-time-ish capability cost; the more trees it serves, the lower its effective
per-tree share.

\paragraph{Inference price and model mix.} The agent slope $\agentcost$ is dominated by routine
monitoring on Haiku; pushing more decisions to Opus, or polling more often, raises it, but it would
have to rise by two orders of magnitude---to $\sim\!\$\num{667}$/shard---before it matched the human
marginal cost. There is enormous headroom: prompt caching, batch settlement of non-urgent checks,
and cheaper routine models all cut $\agentcost$ further. The agent's cost advantage is not fragile.

\paragraph{Span of control.} Better human tooling raises $\sigma_d$ and delays the crossover; it does
not remove it, because the human curve still has positive slope and the agent curve is essentially
flat. No achievable $\sigma_d$ closes a gap that is two orders of magnitude per shard at the margin.

\section{Self-hosting the agent's model: local inference}
\label{sec:local}

Everything so far priced the agent's inference at third-party API rates---pure variable cost, a few
dollars per shard, zero fixed outlay. Dedicating hardware to a \emph{local} model inverts that
structure. A GPU runs whether it serves one shard or three thousand, so local inference is a
\emph{fixed step cost}: the same shape as the human team, one level down the stack. The question is
the same too---at what tree size does the fixed asset amortise below the variable bill?

\paragraph{The commodity tier.} Split the agent's work by volume. Routine monitoring is high-volume
($\sim\!\num{720}$ calls/shard-month) and cognitively cheap---exactly what a small open-weight model
on a commodity inference GPU does well. On an OCI inference shape (an L40S- or A10-class card at
$\sim\$1.50$/GPU-hour list, $\sim\$\num{1095}$/month; $\sim\$0.85$ on a three-year
commit~\cite{oci2026}) running a quantised model at a conservative $\num{5000}$ tokens/s and
\SI{70}{\percent} utilisation, one GPU clears $\sim\!\num{9.2}$ billion tokens/month. The routine
workload is $\num{3.24}$\,M tokens/shard-month, so a single card saturates at $\sim\!\num{2840}$
shards.

\paragraph{Where local beats the API.} Per-shard local cost is the GPU bill divided by the tree it
serves; it falls as the card fills. It crosses the flat API price (\$2.61/shard for routine on Haiku
4.5) at the \emph{break-even tree size}
\[
\shards_{\text{BE}} \;=\; \frac{\text{GPU \$/month}}{\text{API \$/shard}}
\;=\; \frac{\num{1095}}{2.61} \;\approx\; \num{420}\ \text{shards}.
\]
Below \num{420} shards the card sits mostly idle and the API is cheaper; above it, local routine cost
collapses toward $\sim\$0.44$/shard once the GPU saturates---about \SI{17}{\percent} of the API
price, a sixfold cut on the routine tier.

\begin{table}[h]
\centering
\caption*{\textbf{Table 3.} Routine-tier inference: per-shard cost, API vs self-hosted GPU.}
\setlength{\tabcolsep}{7pt}
\renewcommand{\arraystretch}{1.12}
\begin{tabular}{@{}rrrrl@{}}
\toprule
Shards $\shards$ & GPUs & Local \$/shard & API \$/shard & Cheaper \\
\midrule
\num{100}   & 1 & \$10.95 & \$2.61 & API \\
\num{334}   & 1 & \$3.28  & \$2.61 & API \\
\num{420}   & 1 & \$2.61  & \$2.61 & break-even \\
\rowcolor{coolrow}\num{500}   & 1 & \$2.19  & \$2.61 & Local \\
\rowcolor{coolrow}\num{1000}  & 1 & \$1.09  & \$2.61 & Local \\
\rowcolor{coolrow}\num{2500}  & 1 & \$0.44  & \$2.61 & Local \\
\rowcolor{coolrow}\num{12074} & 5 & \$0.45  & \$2.61 & Local \\
\bottomrule
\end{tabular}
\end{table}

\paragraph{Why the frontier tier stays rented.} The hard decisions (split, migrate, triage) run on a
frontier model precisely because quality there is load-bearing, and they are \emph{low} volume
($\sim\!6$/shard-month). Self-hosting frontier quality means a large open-weight model on a
multi-GPU node---an $8\times$H100 box lists near \$\num{58400}/month~\cite{oci2026}---against a
hard-decision API bill of only $\sim\$\num{16800}$/month even at the \num{12074}-shard connector
scale. The dedicated floor is over three times the entire bill it would replace, and open-weight
quality may still trail Opus. The economical architecture is therefore two-tier and mirrors the
human story exactly: \emph{own the commodity high-volume tier, rent the exceptional low-volume one.}

\paragraph{What it changes, and what it doesn't.} Self-hosting does \emph{not} move the agent-versus-DevOps
crossover of \S\ref{sec:cross}. That crossover sits near \num{60} shards, far below the \num{420}-shard
inference break-even, so in the contested region one is on the API regardless; the management decision
is unaffected. What local inference does is \emph{deepen} the agent's lead once the tree is large: the
blended per-shard agent cost falls from \$4.17 to $\sim\$2.01$ at connector scale, and the full-auto
price floor drops from \$\num{80305} to \$\num{54310}/month---a \SI{32}{\percent} cut---because the
single largest variable component, routine inference, has been converted to a nearly-flat asset.

\begin{remark}[Below break-even, non-cost drivers can still decide it]
Cost is not the only axis. Keeping the management model local holds operational telemetry---shard
health, tenant patterns, traffic---off a third-party endpoint, which is a data-residency and
sovereignty argument for the regulated engagements (the VA cohort, Tata, fintech), where the
compliance multipliers in the sector model already show these tenants paying for assurance. It also
shortens the control-loop round trip. Any of these can justify a dedicated GPU \emph{below} the
\num{420}-shard break-even, where it loses on dollars alone. And there is a F1R3FLY-native option the
generic analysis misses: the model can run \emph{as a tenant on the tree it manages}---co-resident,
metered in phlogiston, paying rent like any other residual---so the operator becomes self-hosted
inside the very substrate it operates, rather than a service bolted on beside it.
\end{remark}

\section{Do you even need the model? Algorithmic control and the three-tier hierarchy}
\label{sec:algo}

Everything above assumed that automating the tree means an LLM does the deciding. That assumption is
wrong, and worth correcting. Most shard-tree decisions are \emph{deterministic}: split when
$\alpha>\alpha^\star$, demote a tenant when its rent fund empties, provision a shard when capacity is
hit, remediate when a health metric leaves its band. These are threshold tests and B-tree range
handoffs---exactly the R1--R5 rewrite families of the rent companion---and a conventional control
plane (operators, reconcilers, policy loops) executes them with no model in the loop. The earlier
regimes priced \emph{all} $\sim\!\num{726}$ decisions per shard-month through Haiku or Opus; that
treats the LLM as the control plane rather than as the exception handler above it, and over-prices
inference by roughly twentyfold.

\paragraph{Tier 0: the algorithmic control plane.} A deterministic controller costs a fixed
engineering charge to build and maintain ($\sim\$\num{25000}$/month here) and essentially
\emph{nothing} per shard---the reconcile loop is CPU cycles on nodes you already run. Its fully
automated floor is therefore \emph{flat}: $\sim\$\num{27000}$/month for a tree of \emph{any} size,
from ten shards to twelve thousand. That is the true automation floor, and it sits below every LLM
regime in this note, because it buys no tokens at all.

\paragraph{Tier 1: the model as escalation, priced on the residual.} The LLM's right job is the
fraction of decisions the rules do not settle---genuinely ambiguous triage, novel failure modes,
multi-factor policy calls. If a share $\delta\approx\num{0.95}$ of decisions is deterministic, the
model sees only the residual, and its inference cost falls from \$4.17 to
$a_{\text{resid}}\approx\$0.20$ per shard-month. The LLM is to the rules what the human is to the
LLM: each tier is the exception handler of the tier below it.

\paragraph{Tier 2: the human.} Unchanged---the physical and commercial residual of
\S\ref{sec:hybrid}.

\begin{table}[h]
\centering
\caption*{\textbf{Table 4.} Four ways to run the tree. Floors carry no humans; ``feasible best''
chooses between the two deployable regimes.}
\setlength{\tabcolsep}{5.5pt}
\renewcommand{\arraystretch}{1.12}
\begin{tabular}{@{}rrrrrl@{}}
\toprule
$\shards$ & DevOps & Algo full-auto & Algo $+$ human & Three-tier & Feasible best \\
 & (no auto) & (floor) & (no LLM) & (algo$+$LLM$+$human) & \\
\midrule
\num{50}    & \$\num{40000}   & \$\num{27000} & \cellcolor{coolrow}\$\num{45000}  & \$\num{53010}  & Algo $+$ human \\
\num{150}   & \$\num{100000}  & \$\num{27000} & \cellcolor{coolrow}\$\num{45000}  & \$\num{53030}  & Algo $+$ human \\
\num{250}   & \$\num{180000}  & \$\num{27000} & \$\num{65000}  & \cellcolor{coolrow}\$\num{53050}  & Three-tier \\
\num{1000}  & \$\num{680000}  & \$\num{27000} & \$\num{165000} & \cellcolor{coolrow}\$\num{113200} & Three-tier \\
\num{5000}  & \$\num{3340000} & \$\num{27000} & \$\num{705000} & \cellcolor{coolrow}\$\num{374000} & Three-tier \\
\num{12074} & \$\num{8060000} & \$\num{27000} & \$\num{1645000}& \cellcolor{coolrow}\$\num{855416} & Three-tier \\
\bottomrule
\end{tabular}
\end{table}

\paragraph{What the regimes say.} Among floors, algorithmic full-auto ($\$\num{27000}$ flat)
undercuts the LLM full-auto floor at every scale---no surprise, since it runs no inference. Among
\emph{deployable} options, the ranking flips with scale. Below $\sim\!\num{150}$ shards a rules-only
plane with a single on-call engineer is cheapest: when one human already covers the escalations, the
LLM tier's platform ($\sim\$\num{8000}$/month) and inference are not worth their keep. Above
$\sim\!\num{250}$ shards the three-tier system wins, and the gap widens fast: the model roughly
doubles human reach ($\sigma$ from \num{150} to \num{300}), and at scale that saves \emph{multiple}
\$\num{20000}/month engineers for a few thousand dollars of platform and pennies of residual
inference. The crossover is the point where a no-LLM team must hire its second escalation
engineer. Notably the correctly-scoped three-tier even beats the all-LLM ``hybrid'' of
\S\ref{sec:cross} ($\$\num{855}$k vs $\$\num{898}$k at connector scale), because deterministic
decisions no longer pay for tokens.

\paragraph{The honest conclusion.} The cheapest automation is mostly \emph{rules}, with a thin model
layer on top and a thinner human layer above that. The LLM earns its place precisely in the band of
decisions that are too ambiguous for a deterministic rule yet not physical---and that band, not the
whole workload, is what it should be priced on. None of this moves the automation-versus-DevOps
crossover near \num{60} shards; it lowers the automation floor and corrects the architecture. There
is a F1R3FLY reading too: the control plane is the reducer's own lifecycle---the splits, demotions,
and migrations are rewrites the rho machine already performs---so Tier 0 is less a new system than the
normalisation dynamics given an actuator, with the model invoked only where the rewrite rules leave
genuine policy slack.

\section{One open operational knob}

The model treats provisioning as an amortised monthly episode. In a tree that is actively splitting,
provisioning is bursty---a wave of stand-ups follows every density event---and the agent's value is
highest precisely there, where a human team would be paged into a scramble. The incentive-compatible
refinement is to let the split-overhead $K$ from the rent companion \emph{also} price the management
burst: a split that is cheap to coordinate (large $K$, well-separated join graph) is cheap to
\emph{automate}, while a split that fragments a hot join graph is both consensus-expensive and
attention-expensive. We flag this coupling---management cost as a function of $K$, not a flat
per-shard charge---as the natural next iteration, and leave it as a decision rather than settling it
here.

\bigskip
\noindent\textit{Pieces assembled: three operating regimes and their cost laws (\S\ref{sec:regimes},
\S\ref{sec:models}); a transparent token build-up of the agent's per-shard cost (Table 1); the
crossover at $\sim\!60$ shards with its closed form and its two asymptotes (\S\ref{sec:cross}); the
case for hybrid and the irreducible human residual (\S\ref{sec:hybrid}); a sensitivity sweep over
the levers that actually move it (\S\ref{sec:sens}); and the nested local-vs-API inference crossover
when hardware is dedicated to a self-hosted model (\S\ref{sec:local}); and the pure-algorithmic
control plane with the model recast as a residual escalation tier in a three-tier hierarchy
(\S\ref{sec:algo}). Live model: the
\texttt{Mgmt Assumptions},
\texttt{Agent Cost}, \texttt{Mgmt Crossover}, \texttt{Local Inference}, and
\texttt{Algorithmic vs LLM} tabs of \texttt{f1r3fly\_sharding\_economics.xlsx}.}

\begin{thebibliography}{9}
\addcontentsline{toc}{section}{References}

\bibitem{rentnote}
F1R3FLY Industries.
\newblock Rent and shard splitting: a storage-economics companion to occupancy-driven sharding in F1R3FLY.
\newblock Working note, 2026.

\bibitem{anthropicprice2026}
Anthropic.
\newblock Claude API pricing (Haiku 4.5, Sonnet 4.6, Opus 4.8); prompt-caching and batch discounts.
\newblock \url{https://docs.claude.com/en/docs/about-claude/pricing}, accessed June 2026.

\bibitem{builtin2026}
Built In.
\newblock Site Reliability Engineer and DevOps Engineer salary in US, 2026
(avg.\ total comp \$\num{147000}--\$\num{150000}).
\newblock \url{https://builtin.com/salaries}, accessed June 2026.

\bibitem{kore2026}
KORE1.
\newblock Site Reliability Engineer salary guide 2026 (base \$\num{130000}--\$\num{210000};
staff/principal \$\num{300000}+ all-in).
\newblock \url{https://www.kore1.com/sre-salary-guide-2026/}, accessed June 2026.

\bibitem{oci2026}
Oracle Corporation.
\newblock Oracle Cloud Infrastructure pricing (compute, block volume, networking).
\newblock \url{https://www.oracle.com/cloud/price-list/}, accessed June 2026.

\end{thebibliography}

\end{document}
